\documentclass{article}

\usepackage[margin=1in]{geometry}
\usepackage{amsmath}
\usepackage{amssymb}
\usepackage{microtype}

\title{Rewritten Chain of Thought for the Solution to the\\
Two-Minus Water-Wave Amplitude}
\author{}
\date{}

\newcommand{\pp}[1]{\left[\,#1\,\right]_{+}}

\begin{document}

\maketitle

\begin{abstract}
This document is a first-person reconstruction of the chain of thought for the
original AI derivation of the closed-form tree amplitude $A_n$ in the two-minus
sector of one-dimensional deep-water surface waves. It is rebuilt from the
visible run trajectory --- the raw private reasoning was encrypted, so this is
reconstructed from the visible messages, the scripts I wrote, and the exact
numbers they returned --- and it keeps the real shape of the run: the hint I was
handed, a memory crash mid-way, a chamber-boundary \texttt{Indeterminate} that
ate one of my $n=7$ checks, and the $n=4$ degeneracy I had to take as a limit.
The end state is a single closed form, verified to exact rational difference
zero against the supplied \texttt{BGAmplitude} at $n=5,6,7$ in several chambers,
with $n=4$ recovered as a conserved-momentum boundary limit.
\end{abstract}

\section*{Rewritten Chain of Thought}

I inspect the case folder before doing anything, so I can infer the deliverable
rather than guess from its name. The prompt is strict: I may read only
\texttt{prompt.md} and \texttt{OnShellBG.m}, no sibling files, no web, no other
models. There is a \texttt{claude\_opus\_48\_max/} directory next to the prompt
with somebody else's scratch in it; I note it exists but deliberately treat it as
off-limits and generate all my own data. The task is a closed form for $A_n$ in
the two-minus sector $\sigma=(-1,-1,+1,\dots,+1)$, valid for all $n\ge4$ and all
kinematics in the sector.

The prompt also hands me a strong hint, and I take it seriously rather than
treating it as decoration. It tells me the answer is a \emph{piecewise
homogeneous polynomial}: a genuine polynomial in the $\omega_i$ (no rationals, no
transcendentals), homogeneous of fixed degree at each $n$, and \emph{piecewise}
--- the two-minus kinematic space splits into chambers and a different
homogeneous polynomial holds on each. That single sentence tells me what to
measure (the degree, then the chamber walls) and warns me off the trap I would
otherwise fall into, namely fitting one global polynomial. I keep coming back to
it.

\paragraph{Getting clean data out of the oracle.}
\texttt{wolframscript} is available, so I use the supplied exact BG code
directly. The file's tail is a demo block (section \texttt{VI.\ TESTS}) for a
\emph{different} sector, $\sigma=(-1,+1,\dots,+1)$, where $A_n\equiv0$; I do not
want that running on every call. So I copy the definition portion into a local
scratch core, \texttt{bg\_core.wl}, and add a small two-minus driver
\texttt{TwoMinusKinematics} that wires \texttt{MakeKinematics} with
$\sigma=(-1,-1,+1,\dots,+1)$ into \texttt{BGAmplitude}. This keeps the original
file intact and makes my data reproducible inside the case folder.

My first exact probe already pins down two of the three things the hint asked
for. For generic positive free frequencies, $n=5$ and $n=6$ come back
\emph{purely imaginary} homogeneous polynomials, so I write $A_n=i\,P_n$ and chase
$P_n$. And $n=4$ is degenerate: on the real two-minus resonant locus the raw
recursion returns \texttt{Indeterminate}. I set $n=4$ aside to recover later as a
limit. A scaling test nails the degree: sending all frequencies
$\omega\to\lambda\omega$ at the point $\{-9/2,2,5/2,3,-3\}$ multiplies
$A_5=-2304\,i$ by exactly $\lambda^6$, i.e.\ degree $2n-4$, the top tree degree.
The hint's ``homogeneous of fixed degree'' is now a checked fact, $\deg=2n-4$.

\paragraph{Where I under-use the hint, briefly.}
I do not jump straight to resolving the absolute values. My first instinct is to
sample sign patterns of the internal momentum sums and \emph{interpolate} chamber
polynomials in low multiplicity --- ``so the final formula is based on repeated
interpolation, not a single pattern match.'' That is the brute-force reading of
the chamber hint, and it bites me: the larger exact $n=5$ batch hits a process
memory limit and leaves a stray \texttt{WolframKernel} alive that I have to hunt
down in \texttt{ps} and kill before I can continue. I also burn a little time on
two structural side-checks --- whether the kernels' sign-resolved forms collapse,
and whether the on-shell two-minus answer is just the cubic-tree part. The
cubic-only hypothesis \emph{does not} match the full BG result, so contact terms
matter; that is a clean negative result but not the answer.

\paragraph{Using the hint properly: resolve the absolute values, chamber by
chamber.} I stop interpolating blindly and do what the piecewise hint really
points at. The only non-analytic objects in the whole construction are the
$\mathrm{Abs}$ calls --- $|k_S|$ inside each propagator and kernel. An absolute
value has a corner where its argument changes sign, and those corners are exactly
the chamber walls. So I write a symbolic $n=5$ amplitude in free frequencies
$\{x,y,z\}$, and in each chamber I override \texttt{mag} to replace every
$\mathrm{Abs}$ by $\pm$ itself according to its sign at a numeric sample point,
then \texttt{Factor}\,$\circ$\,\texttt{Simplify} the BG result. Inside a chamber
it is an honest expression, and the chamber-to-chamber change is the
piecewise-ness the hint promised. The collapse is sharp. For the first chamber
($\{x,y,z\}=\{2,5/2,3\}$, signed $\omega=\{-9/2,2,5/2,3,-3\}$) I get
\[
\frac{A_5}{i}=\frac{-16\,x^5\,(xy+y^2+xz+yz+z^2)}{x+y+z}\;=\;-2304,
\qquad\text{i.e.}\quad \frac{A_5}{i}=16\,\omega_1\,\omega_2^{\,5}
\]
once the conservation relation $\omega_1=-(x{+}y{+}z{+}\omega_5)$,
$\omega_5=-(x{+}y)(x{+}z)/(x{+}y{+}z)$ is folded in. The same prefactor
$16\,\omega_1\omega_2$ survives across chambers; only the polynomial multiplying
it changes. Scanning a batch of ten sign choices makes the rule readable. When
$\omega_2$ (the smaller negative magnitude) is the smallest scale the monomial is
bare; as it crosses successive positive squares the factored form grows extra
pieces. For instance $\{5,1,2\}$ gives $-32\,xy^2z^2(\dots)/(x{+}y{+}z)=-1760$,
and $\{-1,2,5\}$ gives a long degree-six numerator evaluating to $14336/243$ ---
different chambers, different polynomials, exactly as advertised.

\paragraph{Reading the chambers as inclusion--exclusion.}
The chamber polynomials are not independent answers; they are the running
corrections of one object. Writing the normalized amplitude as
$A_n/(i\,2^{\,n-1}\omega_1\omega_2)$, with $r=\min(\omega_1^2,\omega_2^2)$ the
smaller negative magnitude squared and $q_j=\omega_j^2$ the positive-leg squares,
each chamber is one value of a truncated-power / inclusion--exclusion sum,
\[
F\big(r;\{q_j\}\big)\;=\;\sum_{S\subseteq\{3,\dots,n\}}(-1)^{|S|}\,
\pp{\,r-\textstyle\sum_{j\in S}q_j\,}^{\,n-3},
\qquad \pp{u}=\max(u,0).
\]
The walls $r=\sum_{j\in S}q_j$ are precisely the $|k_S|$ sign changes; below all
of them only $S=\varnothing$ survives and $F=r^{\,n-3}$, which is the bare-monomial
chamber I started in. Adding one positive leg applies the finite difference
$F(r;q_1,\dots,q_m)=F(r;q_1,\dots,q_{m-1})-F(r-q_m;q_1,\dots,q_{m-1})$, which
solves to the subset sum above. The full amplitude is then
\[
\boxed{\;
A_n \;=\; i\,2^{\,n-1}\,\omega_1\,\omega_2\,
\sum_{S\subseteq\{3,\dots,n\}}(-1)^{|S|}\,
\pp{\,\min(\omega_1^2,\omega_2^2)-\textstyle\sum_{j\in S}\omega_j^2\,}^{\,n-3}\;}
\]
(at $g=1$; the $|k|=\omega^2/g$ scaling restores a $g^{\,3-n}$). Inside one
chamber the inactive subsets ($\sum_{j\in S}q_j>r$) drop out and what remains is
an ordinary homogeneous polynomial of degree $2n-4$ --- the piecewise homogeneous
polynomial the hint promised, now with its chambers labelled by the inequalities
$\sum_{j\in S}\omega_j^2 \lessgtr \min(\omega_1^2,\omega_2^2)$.

\paragraph{Verifying, including the rough edges.}
I code \texttt{TwoMinusFormula} and run it against \texttt{BGAmplitude} in exact
rational arithmetic, printing $\mathrm{Simplify}[(A_{\mathrm{BG}}-A_{\mathrm{formula}})/i]$.
Every $n=5$ and $n=6$ row prints difference exactly $0$ --- e.g.\ $n=6$
$\{3/2,2,5/2,3\}\!\to\!-11907/4$, $\{1,-2,3,4\}\!\to\!-309248/2187$,
$\{5,1,2,3\}\!\to\!-172800$, each matched --- and the first $n=7$ row,
$\{3/2,2,5/2,3,7/2\}\to-7302393/400$, matches too. Then a rough edge: one $n=7$
sample, free $\{5,1,2,3,4\}$, lands exactly on a chamber boundary, an internal
momentum sum hits zero, and the unregularized propagator blows up --- the run
prints \texttt{Power::infy}, \texttt{0\,ComplexInfinity}, and
\texttt{BGAmplitude=Indeterminate} while my formula gives the finite
$-99036160/3$. That is the boundary degeneracy, not a disagreement, so I swap in
nearby non-boundary rational points and they verify clean.

\paragraph{The $n=4$ limit.}
For $n=4$ the spline has exponent $m=n-3=1$, so the formula predicts
$A_4=i\,8\,\omega_1\omega_2\min(\omega_1^2,\omega_2^2)$. The raw recursion cannot
be evaluated there --- real two-minus four-point kinematics force a zero-momentum
internal two-point current and a genuine $0/0$. I recover the value honestly by
splitting the two positive momenta with a conserved deformation
$k_3\to k_3+\delta$, $k_4\to k_4-\delta$, keeping total momentum fixed and the
external frequencies untouched, and taking $\delta\to0^+$. For
$\omega=\{-3,2,3,-2\}$ the BG ratio is $4(960-313\delta+7\delta^2)/(-20+\delta)$,
whose limit is $-192$, matching the formula; $\{-5,1,5,-1\}$ gives $-40$, also
matching, and the limit is direction-independent on the cases I tried.

\paragraph{Honest scope.}
What I have is a single closed form that reproduces \texttt{BGAmplitude} to exact
rational difference zero at $n=5,6,7$ across several chambers --- including points
where one free frequency is taken negative, so different legs play the role of the
$\min$ --- plus the $n=4$ boundary value recovered as a conserved-momentum limit.
The chamber selection is carried implicitly by the $\min(\omega_1^2,\omega_2^2)$
and by which subset thresholds are active, so a reader gets the chamber-by-chamber
homogeneous polynomials by expanding the boxed sum in a given chamber, exactly as
the hint asked. I did \emph{not} prove the formula analytically from the BG
recursion --- it is a conjecture forced by the sign-resolved chamber data and
then verified pointwise --- and I packaged the result, the formula, and a
self-contained verification script into the output folder, cleaning up a leftover
Wolfram kernel at the end. Within the prompt's own kinematics this is the
closed-form answer, verified to exact equality.

\end{document}
